% ---------------------------------------------------------
% TABELA EMENTA 34: Introdução à Administração
% ---------------------------------------------------------
\begin{xltabular}{\linewidth}{|X|p{2.5cm}|p{2.5cm}|}
\hline
\multirow{3}{=}{\textcolor{ifscgreen}{\textbf{Unidade Curricular:}}\\[3pt]\textbf{\large Introdução à Administração}} & \multicolumn{2}{p{\dimexpr5cm+2\tabcolsep+\arrayrulewidth\relax}|}{\textcolor{ifscgreen}{\textbf{Semestre:}} \textbf{1º (Ano 1)}} \\ \cline{2-3}
 & \textcolor{ifscgreen}{\textbf{CH EaD*:}} & \textcolor{ifscgreen}{\textbf{CH Total*:}} \\ \cline{2-3}
 & \textbf{00 h} & \textbf{80 h} \\ \hline
\endhead
\multicolumn{3}{|p{\dimexpr\linewidth-2\tabcolsep-2\arrayrulewidth\relax}|}{\textcolor{ifscgreen}{\textbf{Objetivos:}}} \\ \hline
\multicolumn{3}{|p{\dimexpr\linewidth-2\tabcolsep-2\arrayrulewidth\relax}|}{
\begin{itemize}[leftmargin=*,noitemsep,topsep=2pt]
 \item Compreender o papel da Administração nas organizações e sua evolução histórica, relacionando as principais teorias administrativas às transformações econômicas, tecnológicas e organizacionais da sociedade contemporânea.
 \item Reconhecer o perfil profissional, as competências e os campos de atuação do Técnico em Administração.
 \item Diferenciar os tipos de organizações quanto às suas finalidades, estruturas e formas de gestão, considerando os fatores do ambiente interno e externo que influenciam seu funcionamento.
 \item Identificar as funções administrativas, os níveis organizacionais e a estrutura hierárquica em sua aplicação nas rotinas de trabalho.
 \item Reconhecer tendências contemporâneas que impactam as organizações e a prática administrativa.
 \item Compreender as etapas básicas do processo decisório na identificação de problemas em situações organizacionais.
\end{itemize}
} \\ \hline
\multicolumn{3}{|p{\dimexpr\linewidth-2\tabcolsep-2\arrayrulewidth\relax}|}{\textcolor{ifscgreen}{\textbf{Conteúdos:}}} \\ \hline
\multicolumn{3}{|p{\dimexpr\linewidth-2\tabcolsep-2\arrayrulewidth\relax}|}{
\begin{itemize}[leftmargin=*,noitemsep,topsep=2pt]
 \item Conceito, importância e campo de atuação da Administração.
 \item Perfil profissional, competências e áreas de atuação do Técnico em Administração.
 \item Organizações: conceitos, tipos, objetivos, recursos e classificação.
 \item Ambiente organizacional: fatores internos e externos.
 \item Estrutura organizacional, níveis hierárquicos e papéis gerenciais.
 \item Evolução das teorias administrativas.
 \item Funções administrativas: planejamento, organização, direção e controle.
 \item Tendências contemporâneas em administração.
 \item Processo decisório nas organizações: etapas básicas e sua relação com as funções administrativas.
\end{itemize}
} \\ \hline
\multicolumn{3}{|p{\dimexpr\linewidth-2\tabcolsep-2\arrayrulewidth\relax}|}{\textcolor{ifscgreen}{\textbf{Estratégias de Ensino e Aprendizagem:}}} \\ \hline
\multicolumn{3}{|p{\dimexpr\linewidth-2\tabcolsep-2\arrayrulewidth\relax}|}{- **Metodologia:** Desenvolvida por meio de metodologias ativas de aprendizagem, articuladas a aulas expositivas dialogadas, aproximando os estudantes da realidade das organizações e do mundo do trabalho, com análise de situações reais, observação de diferentes contextos organizacionais e resolução de problemas compatíveis com o nível de formação. Além da sala de aula, são utilizados espaços do câmpus como biblioteca, sala multidisciplinar, centro multiuso e laboratório de informática. Atividades: estudos de caso, exercícios, provas, simulações, dinâmicas de grupos, jogos, seminários, debates, análise de filmes, esquetes teatrais, relatórios, atividades reflexivas, autoavaliações, rodas de conversa com profissionais da área e visitas técnicas. Por ser unidade basilar da formação profissional, dialoga com as demais unidades da formação profissional e pode articular-se com Sociologia, Sociedade e Trabalho, Língua Portuguesa, História e Geografia.
- **Avaliação:** Contínua, processual e formativa, acompanhando o desenvolvimento dos estudantes ao longo de toda a unidade curricular, considerando participação, envolvimento nas atividades propostas e atitudes em sala de aula.} \\ \hline
\multicolumn{3}{|p{\dimexpr\linewidth-2\tabcolsep-2\arrayrulewidth\relax}|}{\textcolor{ifscgreen}{\textbf{Bibliografia Básica:}}} \\ \hline
\multicolumn{3}{|p{\dimexpr\linewidth-2\tabcolsep-2\arrayrulewidth\relax}|}{1. DIAS, Reinaldo; ZAVAGLIA, Tércia; CASSAR, Maurício. \textbf{Introdução à administração: da competitividade à sustentabilidade}. 3. ed. rev. Campinas, SP: Alínea, 2013.
2. MAXIMIANO, Antonio Cesar Amaru. \textbf{Introdução à administração}. 8. ed. rev. e ampl. São Paulo: Atlas, 2011.} \\ \hline
\multicolumn{3}{|p{\dimexpr\linewidth-2\tabcolsep-2\arrayrulewidth\relax}|}{\textcolor{ifscgreen}{\textbf{Bibliografia Complementar:}}} \\ \hline
\multicolumn{3}{|p{\dimexpr\linewidth-2\tabcolsep-2\arrayrulewidth\relax}|}{1. CHIAVENATO, Idalberto. \textbf{Administração nos novos tempos}. 2. ed. rev. e atual. Rio de Janeiro: Elsevier, 2010.
2. LACOMBE, Francisco José Masset; HEILBORN, Gilberto Luiz José. \textbf{Administração: princípios e tendências}. 3. ed. São Paulo: Saraiva, 2015.
3. NASCIMENTO, Edson Ronaldo. \textbf{Gestão Pública}. São Paulo: Saraiva, 2010.} \\ \hline
\end{xltabular}

\vspace{0.4cm}


% ---------------------------------------------------------
% TABELA EMENTA 35: Sociedade e Trabalho
% ---------------------------------------------------------
\begin{xltabular}{\linewidth}{|X|p{2.5cm}|p{2.5cm}|}
\hline
\multirow{3}{=}{\textcolor{ifscgreen}{\textbf{Unidade Curricular:}}\\[3pt]\textbf{\large Sociedade e Trabalho}} & \multicolumn{2}{p{\dimexpr5cm+2\tabcolsep+\arrayrulewidth\relax}|}{\textcolor{ifscgreen}{\textbf{Semestre:}} \textbf{Não especificado no documento fonte}} \\ \cline{2-3}
 & \textcolor{ifscgreen}{\textbf{CH EaD*:}} & \textcolor{ifscgreen}{\textbf{CH Total*:}} \\ \cline{2-3}
 & \textbf{00 h} & \textbf{40 h} \\ \hline
\endhead
\multicolumn{3}{|p{\dimexpr\linewidth-2\tabcolsep-2\arrayrulewidth\relax}|}{\textcolor{ifscgreen}{\textbf{Objetivos:}}} \\ \hline
\multicolumn{3}{|p{\dimexpr\linewidth-2\tabcolsep-2\arrayrulewidth\relax}|}{
\begin{itemize}[leftmargin=*,noitemsep,topsep=2pt]
 \item Compreender as transformações que ocorrem nas sociedades humanas e sua relação com o trabalho.
 \item Evidenciar a relação entre as questões individuais e as questões sociais.
 \item Formular questionamentos que permitam alcançar um conhecimento mais preciso da sociedade e do trabalho, possibilitando a construção de uma postura crítica em relação às vivências que nos condicionam e limitam.
 \item Entender o processo de constituição, consolidação e desenvolvimento das sociedades modernas e a importância do trabalho na dinâmica de funcionamento das sociedades.
\end{itemize}
} \\ \hline
\multicolumn{3}{|p{\dimexpr\linewidth-2\tabcolsep-2\arrayrulewidth\relax}|}{\textcolor{ifscgreen}{\textbf{Conteúdos:}}} \\ \hline
\multicolumn{3}{|p{\dimexpr\linewidth-2\tabcolsep-2\arrayrulewidth\relax}|}{
\begin{itemize}[leftmargin=*,noitemsep,topsep=2pt]
 \item O que é Trabalho.
 \item História do Trabalho.
 \item Revolução Industrial e Trabalho assalariado.
 \item Trabalho nas sociedades modernas.
 \item Formas de organização do mundo do Trabalho.
 \item Movimentos sociais e Trabalho.
 \item Trabalho e Tempo Livre.
 \item Trabalho e Lazer.
 \item Trabalho, Cidadania e Direitos Humanos.
 \item Novas modalidades de trabalho.
\end{itemize}
} \\ \hline
\multicolumn{3}{|p{\dimexpr\linewidth-2\tabcolsep-2\arrayrulewidth\relax}|}{\textcolor{ifscgreen}{\textbf{Estratégias de Ensino e Aprendizagem:}}} \\ \hline
\multicolumn{3}{|p{\dimexpr\linewidth-2\tabcolsep-2\arrayrulewidth\relax}|}{- **Metodologia:** Desenvolvida por meio de leitura, análise, pesquisas, discussão e exposição de textos e imagens em sala de aula, visando o exercício do debate e da reflexão crítica sobre os temas e conceitos estudados; aulas expositivas e dialogadas; recursos didático-pedagógicos como filmes, seminários, documentários e entrevistas; estímulo à autonomia investigativa e socialização de temas relacionados ao programa curricular. Recursos didáticos: textos (livros, apostilas, artigos, estudos de casos), quadro e pincel, recursos audiovisuais (filmes, séries, música, pesquisas, arte) e equipamentos de informática (retroprojetor, computador, internet).
- **Avaliação:** Não especificado no documento fonte.} \\ \hline
\multicolumn{3}{|p{\dimexpr\linewidth-2\tabcolsep-2\arrayrulewidth\relax}|}{\textcolor{ifscgreen}{\textbf{Bibliografia Básica:}}} \\ \hline
\multicolumn{3}{|p{\dimexpr\linewidth-2\tabcolsep-2\arrayrulewidth\relax}|}{1. Livro Didático fornecido pelo Fundo Nacional de Desenvolvimento da Educação (FNDE).
2. OLIVEIRA, Pérsio Santos de. \textbf{Introdução à Sociologia: Série Brasil}. São Paulo: Editora Ática, 2011.} \\ \hline
\multicolumn{3}{|p{\dimexpr\linewidth-2\tabcolsep-2\arrayrulewidth\relax}|}{\textcolor{ifscgreen}{\textbf{Bibliografia Complementar:}}} \\ \hline
\multicolumn{3}{|p{\dimexpr\linewidth-2\tabcolsep-2\arrayrulewidth\relax}|}{1. COSTA, Maria Cristina. \textbf{Sociologia: introdução à ciência da sociedade}. 5. ed. São Paulo: Editora Moderna, 2016.
2. OLIVEIRA, Luiz Fernandes de; COSTA, Ricardo Cesar Rocha da. \textbf{Sociologia para jovens do século XXI}. Rio de Janeiro: Imperial Novo Milênio, 2013.
3. HELENA, Bomeny et al. \textbf{Tempos Modernos, tempos de sociologia}. 4. ed. São Paulo: Ed. do Brasil, 2016.} \\ \hline
\end{xltabular}

\vspace{0.4cm}


% ---------------------------------------------------------
% TABELA EMENTA 36: Matemática para Administração
% ---------------------------------------------------------
\begin{xltabular}{\linewidth}{|X|p{2.5cm}|p{2.5cm}|}
\hline
\multirow{3}{=}{\textcolor{ifscgreen}{\textbf{Unidade Curricular:}}\\[3pt]\textbf{\large Matemática para Administração}} & \multicolumn{2}{p{\dimexpr5cm+2\tabcolsep+\arrayrulewidth\relax}|}{\textcolor{ifscgreen}{\textbf{Semestre:}} \textbf{3º (Ano 2)}} \\ \cline{2-3}
 & \textcolor{ifscgreen}{\textbf{CH EaD*:}} & \textcolor{ifscgreen}{\textbf{CH Total*:}} \\ \cline{2-3}
 & \textbf{00 h} & \textbf{40 h} \\ \hline
\endhead
\multicolumn{3}{|p{\dimexpr\linewidth-2\tabcolsep-2\arrayrulewidth\relax}|}{\textcolor{ifscgreen}{\textbf{Objetivos:}}} \\ \hline
\multicolumn{3}{|p{\dimexpr\linewidth-2\tabcolsep-2\arrayrulewidth\relax}|}{
\begin{itemize}[leftmargin=*,noitemsep,topsep=2pt]
 \item Empregar conceitos de matemática básica para dar suporte à execução de rotinas operacionais.
 \item Efetuar operações de matemática financeira aplicadas à administração.
 \item Aplicar métodos estatísticos básicos para auxiliar na tomada de decisões gerenciais.
\end{itemize}
} \\ \hline
\multicolumn{3}{|p{\dimexpr\linewidth-2\tabcolsep-2\arrayrulewidth\relax}|}{\textcolor{ifscgreen}{\textbf{Conteúdos:}}} \\ \hline
\multicolumn{3}{|p{\dimexpr\linewidth-2\tabcolsep-2\arrayrulewidth\relax}|}{
\begin{itemize}[leftmargin=*,noitemsep,topsep=2pt]
 \item Juros simples.
 \item Desconto simples.
 \item Juros compostos.
 \item Inflação e correção.
 \item Sistemas de amortização.
 \item Medidas de posição: média, moda, mediana, quartis.
 \item Medidas de dispersão: desvio médio, variância, desvio padrão.
\end{itemize}
} \\ \hline
\multicolumn{3}{|p{\dimexpr\linewidth-2\tabcolsep-2\arrayrulewidth\relax}|}{\textcolor{ifscgreen}{\textbf{Estratégias de Ensino e Aprendizagem:}}} \\ \hline
\multicolumn{3}{|p{\dimexpr\linewidth-2\tabcolsep-2\arrayrulewidth\relax}|}{- **Metodologia:** Aulas expositivas e dialogadas com o objetivo de relacionar os conhecimentos e sua aplicabilidade no mundo produtivo, instigando o aluno na resolução de cases e interpretação de dados. São trabalhadas competências como: realizar cálculos de matemática financeira aplicados à administração; utilizar planilha eletrônica para operações financeiras; transcrever mensagens matemáticas da linguagem corrente para linguagem simbólica (equações, gráficos, diagramas, fórmulas, tabelas etc.) e vice-versa; elaborar e interpretar gráficos e tabelas estatísticas; definir e calcular medidas de posição e de dispersão; utilizar softwares para cálculos estatísticos; e aplicar os conteúdos apresentados em situações do cotidiano para resolução de problemas.
- **Avaliação:** Consistirá de resolução de exercícios.} \\ \hline
\multicolumn{3}{|p{\dimexpr\linewidth-2\tabcolsep-2\arrayrulewidth\relax}|}{\textcolor{ifscgreen}{\textbf{Bibliografia Básica:}}} \\ \hline
\multicolumn{3}{|p{\dimexpr\linewidth-2\tabcolsep-2\arrayrulewidth\relax}|}{1. CRESPO, A. A. \textbf{Estatística fácil}. 19. ed. São Paulo: Saraiva, 2009.
2. DANTE, L. R. \textbf{Matemática contexto e aplicações}. 4. ed. São Paulo: Ática, 2011.} \\ \hline
\multicolumn{3}{|p{\dimexpr\linewidth-2\tabcolsep-2\arrayrulewidth\relax}|}{\textcolor{ifscgreen}{\textbf{Bibliografia Complementar:}}} \\ \hline
\multicolumn{3}{|p{\dimexpr\linewidth-2\tabcolsep-2\arrayrulewidth\relax}|}{1. BARBETTA, P. A. \textbf{Estatística aplicada às ciências sociais}. Florianópolis: Edufsc, 2010.
2. BRUNI, Adriano Leal; FAMÁ, Rubens. \textbf{Matemática financeira com HP 12C e Excel}. São Paulo: Atlas, 2008.
3. BUENO, Fabrício. \textbf{Estatística para processos produtivos}. Florianópolis: Visual Books, 2010.
4. BUENO, Fabrício. \textbf{Estaticópolis: um jeito novo de aprender Estatística}. São Paulo: Amazon, 2015.} \\ \hline
\end{xltabular}

\vspace{0.4cm}

